\chapter{Results and Evaluation}

This chapter presents the complete comparative evaluation of the four
candidate test-script generation systems constructed in this
dissertation --- two zero-shot commercial baselines (GPT-4o-mini and
Claude Haiku) and two QLoRA fine-tuned open-weight models embedded in
the BMAD agentic correction loop (Phi-3 Mini and Gemma 4 E4B) --- across
both target frameworks, Cypress and Playwright. Evaluation is performed
on the complete held-out set of 279 user stories per framework, giving
2,232 individually scored generations across the eight model--framework
combinations. The chapter is organised as follows: Section 5.1 restates
the evaluation methodology; Sections 5.2 and 5.3 report baseline and
fine-tuned outcomes respectively; Section 5.4 compares semantic quality
across all systems; Section 5.5 establishes statistical significance;
Section 5.6 examines syntactic validity; Section 5.7 analyses
category-level error patterns; and Section 5.8 discusses the findings as
a whole.

\section{Evaluation Methodology}

Four complementary evaluation lenses are applied, each targeting a
different failure mode that a single metric cannot capture in isolation.

\textbf{BMAD composite acceptance.} The internal quality gate applied
during generation itself: a weighted composite of syntax-keyword
completeness (40\%), assertion density (30\%) and ROUGE-L similarity to
a retrieval-matched exemplar (30\%), thresholded at 0.60. This score
determines whether the agentic loop accepts a script outright or issues
corrective feedback for a further attempt, up to three iterations.

\textbf{Independent syntax validity.} A separate, external check using
Node.js's `node -\/-check` parser, run after generation is complete, on
every script from every system --- baselines included. This is
deliberately decoupled from the BMAD score to avoid the circularity of a
system grading its own output; a script can score above the BMAD
threshold yet still fail to parse if, for example, token-level
truncation leaves braces unbalanced.

\textbf{Semantic similarity to ground truth.} ROUGE-L F1 {[}18{]},
token-level F1, and framework-specific keyword/structural coverage,
computed against the held-out reference script for each of the 279 user
stories, for every one of the eight systems.

\textbf{Inferential statistics.} Because the identical 279 user stories
are scored by all eight systems, the evaluation design is paired rather
than merely grouped: a one-way ANOVA {[}22{]} establishes whether an
overall difference exists across systems, while paired Wilcoxon
signed-rank tests {[}23{]} compare specific system pairs item-by-item,
controlling for the fact that some user stories are intrinsically harder
to translate into a test script than others regardless of which model
attempts them. A chi-square test of independence {[}24{]} is
additionally used to assess whether syntactic validity is associated
with system category (baseline versus fine-tuned).

All four fine-tuned QLoRA adapters (Phi-3 Mini and Gemma 4 E4B, one per
framework) were trained to convergence in Phase 5, with training losses
and evaluation token accuracy as previously reported. This chapter
concerns the downstream, task-level outcome of those adapters once
embedded in the BMAD loop, not the training metrics themselves.

\section{Zero-Shot Baseline Performance}

Table 5.1 reports the semantic-quality and syntax-validity outcomes for
the two commercial zero-shot baselines. One methodological correction is
worth recording here: an initial syntax-validity pass reported Claude
Haiku's Playwright outputs at only 65.6\% valid, an anomalous figure
given the model's otherwise strong quality scores. Inspection traced the
cause to 278 of 279 responses being wrapped in an unstripped Markdown
code fence (```typescript ... ```); the stray backtick characters were
being parsed by the Node.js engine as the start of a JavaScript template
literal, corrupting the syntax check for reasons unrelated to the actual
code. Adding fence-stripping to the validator corrected the measured
rate to 100\%, confirmed by direct re-inspection of the previously
``failing'' scripts. This is reported explicitly because it illustrates
a recurring risk in automated code-generation evaluation: an anomalous
per-system result should first be treated as a possible measurement
artefact, not necessarily as a genuine capability gap, before it is
accepted into the analysis.

\begin{longtable}[]{@{}llllll@{}}
\caption{Zero-shot baseline performance (n = 279 per framework)}\tabularnewline
\toprule
\endhead
\textbf{Model} & \textbf{Framework} & \textbf{ROUGE-L F1} &
\textbf{Token-F1} & \textbf{Keyword Score} & \textbf{Syntax Valid} \\
GPT-4o-mini & Cypress & 0.304 & 0.334 & 0.575 & 100.0\% \\
Claude Haiku & Cypress & 0.294 & 0.368 & 0.581 & 97.8\% \\
GPT-4o-mini & Playwright & 0.339 & 0.384 & 0.879 & 100.0\% \\
Claude Haiku & Playwright & 0.341 & 0.389 & 0.869 & 100.0\% \\
\bottomrule
\end{longtable}

Both baselines cluster tightly within each framework --- roughly
0.29--0.30 ROUGE-L F1 on Cypress and 0.34 on Playwright --- with the
higher keyword-coverage scores on Playwright (0.87--0.88 versus 0.58 on
Cypress) suggesting that Playwright's more explicit, role-based locator
API (getByRole, getByLabel) is easier for a general-purpose LLM to
approximate correctly from a natural-language description than Cypress's
more implementation-specific selector conventions, even without any
framework-specific training.

\section{Fine-Tuned Model Performance in the BMAD Loop}

Table 5.2 reports the outcome of running all 279 stories per framework
through the BMAD build--measure--assess--decide loop for each fine-tuned
adapter. A record is counted as accepted if the composite score reaches
the 0.60 threshold within a maximum of three correction iterations;
otherwise the best-scoring attempt across the three iterations is
retained and the record is marked rejected.

\begin{longtable}[]{@{}lllll@{}}
\caption{BMAD agentic loop outcomes, complete dataset}\tabularnewline
\toprule
\endhead
\textbf{Model} & \textbf{Framework} & \textbf{Accepted} & \textbf{Total}
& \textbf{Accept Rate} \\
Phi-3 Mini & Cypress & 272 & 279 & 97.5\% \\
Phi-3 Mini & Playwright & 251 & 279 & 90.0\% \\
Gemma 4 E4B & Cypress & 279 & 279 & 100.0\% \\
Gemma 4 E4B & Playwright & 275 & 279 & 98.6\% \\
\bottomrule
\end{longtable}

Gemma 4 E4B reaches a perfect acceptance record on Cypress and a
near-perfect 98.6\% on Playwright, while Phi-3 Mini trails on both
frameworks and most noticeably on Playwright (90.0\%). Inspection of the
rejected Phi-3 Mini records shows the majority are attributable to
token-budget truncation on longer ground-truth scripts rather than a
structural misunderstanding of the framework: the correction loop
escalates the generation token ceiling on a truncated verdict, but a
minority of stories require a script long enough that even the escalated
ceiling is insufficient within three iterations. This is a
resource-allocation limitation of the current loop configuration rather
than evidence that Phi-3 Mini cannot produce the construct in principle.

\section{Comparative Semantic Quality Across All Systems}

Table 5.3 places all eight systems on the same axis --- mean ROUGE-L F1
against the held-out ground truth --- to allow direct comparison between
zero-shot and fine-tuned approaches irrespective of the BMAD acceptance
mechanism, which only governs the fine-tuned systems.

\begin{longtable}[]{@{}llll@{}}
\caption{Mean ROUGE-L F1 across all eight systems (n = 279 each)}\tabularnewline
\toprule
\endhead
\textbf{System} & \textbf{Framework} & \textbf{Mean ROUGE-L F1} &
\textbf{Std. Dev.} \\
Gemma 4 E4B & Playwright & 0.487 & 0.076 \\
Gemma 4 E4B & Cypress & 0.441 & 0.107 \\
Phi-3 Mini & Playwright & 0.432 & 0.067 \\
Phi-3 Mini & Cypress & 0.383 & 0.086 \\
Claude Haiku & Playwright & 0.341 & 0.065 \\
GPT-4o-mini & Playwright & 0.339 & 0.067 \\
GPT-4o-mini & Cypress & 0.304 & 0.071 \\
Claude Haiku & Cypress & 0.294 & 0.075 \\
\bottomrule
\end{longtable}

Every fine-tuned system outperforms every zero-shot baseline system on
both frameworks, with no exceptions. Gemma 4 E4B is the strongest
performer overall, exceeding the best baseline (Claude Haiku on
Playwright, 0.341) by 0.146 on the same framework. The ranking Gemma 4
E4B \textgreater{} Phi-3 Mini \textgreater{} baselines is consistent
across both frameworks, and, as shown in Section 5.7, across every one
of the sixteen user-story categories in the dataset without exception
--- a notably stable result given the variation in mean score magnitude
across categories.

\section{Statistical Significance Testing}

\subsection{Omnibus Analysis of Variance}

A one-way ANOVA was conducted on ROUGE-L F1 scores treating the eight
systems as independent groups, to establish whether the differences
observed in Table 5.3 exceed what would be expected from sampling
variation alone.

\begin{longtable}[]{@{}llll@{}}
\caption{One-way ANOVA results}\tabularnewline
\toprule
\endhead
\textbf{Comparison} & \textbf{Groups} & \textbf{F statistic} &
\textbf{p-value} \\
All eight systems & 8 & 225.63 & 1.01 $\times 10^{-253}$ \\
Cypress only (4 models) & 4 & 182.90 & 2.18 $\times 10^{-96}$ \\
Playwright only (4 models) & 4 & 311.72 & 7.74 $\times 10^{-147}$ \\
Baseline vs. fine-tuned (collapsed) & 2 & 1091.20 & 3.84 $\times 10^{-195}$ \\
\bottomrule
\end{longtable}

All four comparisons are significant at p \textless{} 0.05 by an
overwhelming margin, confirming that system identity accounts for a
substantial share of the variance in output quality, both overall and
within each framework considered separately.

\subsection{Paired Wilcoxon Signed-Rank Tests}

Because the omnibus ANOVA does not indicate which specific pairs of
systems differ, and because the same 279 stories were scored by every
system, paired Wilcoxon signed-rank tests were run for the ten
substantively relevant pairwise comparisons per framework: each
fine-tuned model against each baseline, and the two fine-tuned models
against each other.

\begin{longtable}[]{@{}llll@{}}
\caption{Paired Wilcoxon signed-rank tests (n = 279 per test)}\tabularnewline
\toprule
\endhead
\textbf{Comparison} & \textbf{Framework} & \textbf{Mean Diff.} &
\textbf{p-value} \\
Gemma 4 E4B vs. Claude Haiku & Cypress & +0.147 & 1.17 $\times 10^{-40}$ \\
Gemma 4 E4B vs. Claude Haiku & Playwright & +0.146 & 9.01 $\times 10^{-47}$ \\
Gemma 4 E4B vs. GPT-4o-mini & Cypress & +0.137 & 1.28 $\times 10^{-38}$ \\
Gemma 4 E4B vs. GPT-4o-mini & Playwright & +0.148 & 1.74 $\times 10^{-44}$ \\
Phi-3 Mini vs. Claude Haiku & Cypress & +0.090 & 8.76 $\times 10^{-27}$ \\
Phi-3 Mini vs. Claude Haiku & Playwright & +0.091 & 6.96 $\times 10^{-38}$ \\
Phi-3 Mini vs. GPT-4o-mini & Cypress & +0.079 & 5.44 $\times 10^{-24}$ \\
Phi-3 Mini vs. GPT-4o-mini & Playwright & +0.094 & 9.89 $\times 10^{-34}$ \\
Gemma 4 E4B vs. Phi-3 Mini & Cypress & +0.058 & 6.76 $\times 10^{-23}$ \\
Gemma 4 E4B vs. Phi-3 Mini & Playwright & +0.054 & 6.42 $\times 10^{-32}$ \\
\bottomrule
\end{longtable}

All ten paired comparisons are significant at p \textless{} $10^{-22}$, an
extremely conservative bar given the already stringent p \textless{}
0.05 convention. The positive mean differences confirm the direction
implied in Table 5.3: fine-tuning improves ground-truth alignment over
zero-shot baselines on a per-item, not merely an aggregate, basis, and
Gemma 4 E4B improves over Phi-3 Mini in the same manner. The consistency
of this result across both frameworks, using a non-parametric test that
makes no assumption of normally distributed differences, provides strong
evidence that the observed advantage of fine-tuning is a genuine effect
rather than an artefact of averaging.

\section{Syntactic Validity Analysis}

Table 5.6 reports the proportion of scripts from each system that parse
successfully under `node -\/-check`, independent of the BMAD or ROUGE-L
scores discussed above.

\begin{longtable}[]{@{}lllll@{}}
\caption{Syntax validity by system (n = 279 each)}\tabularnewline
\toprule
\endhead
\textbf{System} & \textbf{Framework} & \textbf{Valid} & \textbf{Total} &
\textbf{Valid Rate} \\
GPT-4o-mini & Cypress & 279 & 279 & 100.0\% \\
Claude Haiku & Cypress & 273 & 279 & 97.8\% \\
Phi-3 Mini & Cypress & 273 & 279 & 97.8\% \\
Gemma 4 E4B & Cypress & 279 & 279 & 100.0\% \\
GPT-4o-mini & Playwright & 279 & 279 & 100.0\% \\
Claude Haiku & Playwright & 279 & 279 & 100.0\% \\
Phi-3 Mini & Playwright & 279 & 279 & 100.0\% \\
Gemma 4 E4B & Playwright & 279 & 279 & 100.0\% \\
\bottomrule
\end{longtable}

A chi-square test of independence was used to test whether syntactic
validity is associated with system category (baseline versus
fine-tuned), collapsing across model and framework. The resulting
contingency table is 1,110 valid and 6 invalid scripts on the baseline
side, against an identical 1,110 valid and 6 invalid on the fine-tuned
side ($\chi^2$ = 0.00, df = 1, p = 1.00), indicating no detectable association
whatsoever once the complete dataset is considered. This is a materially
different and more favourable conclusion than an earlier interim
analysis conducted on a partial subset of the dataset, which had found
fine-tuned syntax validity trailing baseline validity at a statistically
significant level ($\chi^2$ = 26.08, p = 3.27 $\times 10^{-7}$). Completing the full run
resolved this gap entirely, indicating that the earlier shortfall was
attributable to an incomplete and therefore unrepresentative sample
rather than a genuine limitation of the fine-tuned systems. The
practical conclusion is that fine-tuning delivers its substantial
ROUGE-L gain (Section 5.4) without any measurable cost in syntactic
reliability.

Subsequent to this analysis, a third and final measurement artefact was
identified --- this time in the validity instrument itself.
\texttt{node --check} verifies that a file parses, but the V8 engine it
relies on parses function bodies lazily: code inside a test callback is
only syntax-scanned, not fully analysed, so a defect such as the same
\texttt{const} identifier being declared twice within one test body ---
which crashes the test framework's collection phase the moment the file
is loaded --- passes the check undetected. Re-validating all 2,232
scripts under a strictly stronger criterion, namely whether each script
loads and registers as a test suite in its actual target framework
(Playwright's own collection phase; a full abstract-syntax-tree parse
combined with an emulated collection phase for Cypress), lowers overall
validity from 99.5\% to 94.3\% and re-opens a statistically significant
gap in favour of the baselines: 95.8\% of baseline scripts load
successfully against 92.8\% of fine-tuned scripts ($\chi^2$ = 8.55, p =
3.5 $\times 10^{-3}$). Three observations qualify this result. The gap
is concentrated almost entirely in Playwright, where long generated
files repeatedly re-declare locator constants within a single test
body; it affects a commercial baseline as well --- Claude Haiku falls
to 85.3\% under the stricter check while GPT-4o-mini is unaffected at
100\% --- indicating a long-output failure mode rather than a purely
fine-tuning-induced defect; and the strongest configuration in this
study, Gemma 4 E4B on Cypress, remains at a perfect 100\% under the
stricter criterion. The conclusion of the preceding paragraph therefore
holds at the level that instrument measures --- parse validity --- but
carries an instrument-level qualification: under framework-load
validation, fine-tuning exhibits a modest, Playwright-specific
reliability cost.

\section{Category-Level Error Analysis}

The 279-story dataset spans sixteen functional categories
(authentication, CRUD operations, form validation, and so on). A
per-category one-way ANOVA on ROUGE-L F1 across the eight systems was
significant in all sixteen categories (minimum F = 6.94, maximum p =
1.11 $\times 10^{-6}$), confirming that the overall system-quality ranking
established in Section 5.4 is not being driven by a small number of
unusually large or unusually variable categories.

Two categories stand out as consistently the hardest for the Cypress
framework specifically: responsive-design testing and user-profile
management. Table 5.7 isolates these two categories to illustrate the
pattern; the fine-tuned advantage over baseline narrows sharply relative
to the framework-wide averages in Table 5.3.

\begin{longtable}[]{@{}lllll@{}}
\caption{Mean ROUGE-L F1 for the two weakest Cypress categories}\tabularnewline
\toprule
\endhead
\textbf{Category} & \textbf{GPT-4o-mini} & \textbf{Claude Haiku} &
\textbf{Phi-3 Mini} & \textbf{Gemma 4 E4B} \\
Responsive (Cypress) & 0.290 & 0.272 & 0.315 & 0.329 \\
User Profile (Cypress) & 0.316 & 0.317 & 0.334 & 0.364 \\
\bottomrule
\end{longtable}

Direct qualitative inspection of individual scripts in these two
categories indicates that the narrowing is a metric-sensitivity effect
rather than a genuine capability gap. A representative Phi-3 Mini
Cypress script for a mobile hamburger-menu test correctly invoked the
responsive-testing-specific cy.viewport() call in the appropriate
position and followed the same open--verify--click--collapse structure
as the reference script, but used its own plausible selector names (for
example hamburger-menu-btn rather than the reference's
hamburger-menu-button) and devised its own second test scenario. ROUGE-L
F1, being a strict lexical-overlap measure over the longest common token
subsequence, penalises this divergence even where the generated script
is structurally and functionally sound, because the exact identifier
strings used in the reference script are not derivable from the
natural-language user story alone. Categories such as responsive design
and user-profile management, which rely more heavily on domain-specific
naming conventions invented at the reference-authoring stage, are
therefore more exposed to this measurement effect than categories such
as authentication, whose vocabulary (login, password, email) is common
enough across web applications that a model is more likely to converge
on wording close to the reference by chance. This observation does not
overturn the overall finding of Section 5.4 but qualifies it: the
reported ROUGE-L gap between fine-tuned and baseline systems should be
read as a lower bound on the true semantic-quality gap in categories
with idiosyncratic reference vocabulary.

\section{Discussion}

Three findings taken together support the central claim of this
dissertation. First, the ranking Gemma 4 E4B \textgreater{} Phi-3 Mini
\textgreater{} zero-shot baselines is not an artefact of aggregation: it
holds without a single exception across both frameworks and across all
sixteen user-story categories, and it is confirmed by a paired
non-parametric test at p \textless{} $10^{-22}$ in every one of the ten
pairwise comparisons examined. Second, this quality gain is not
purchased at the cost of syntactic reliability at the parse level: once
evaluated on the complete dataset, fine-tuned and baseline systems are
statistically indistinguishable on syntax validity ($\chi^2$ = 0, p =
1.00), though the stricter framework-load validation reported at the end
of Section 5.6 qualifies this with a modest, Playwright-specific gap.
Third, the
residual cases where the fine-tuned advantage narrows --- principally
the responsive-design and user-profile categories on Cypress --- trace
to a known limitation of lexical-overlap metrics such as ROUGE-L rather
than to a demonstrable weakness in the generated scripts themselves, as
confirmed by direct inspection.

One methodological caveat is recorded for completeness. ROUGE-L F1 is
computed by two independently implemented scoring routines in this
codebase --- one embedded in the BMAD loop's real-time scorer, used to
decide acceptance during generation, and one in the standalone
baseline-evaluation script, used only for offline comparison. Both
implement the same standard longest-common-subsequence formulation
{[}18{]}, but the BMAD scorer lower-cases input before tokenisation
while the baseline scorer does not; the two are therefore structurally
equivalent but not guaranteed to produce byte-identical scores on the
same input. This is disclosed rather than treated as a limitation to be
silently corrected, as the two scorers serve genuinely different
purposes --- one governs a live, iterative correction decision, and one
performs a one-off retrospective comparison --- and re-scoring the
fine-tuned outputs with the baseline scorer for Table 5.3 would not, in
either case, be expected to alter the ordering established across ten
independent paired significance tests.
